% Generated draft: requires mathematical and source review.
\documentclass[11pt]{article}
\usepackage[margin=25mm]{geometry}
\usepackage{fontspec}
\setmainfont{Latin Modern Roman}
\usepackage{amsmath,amssymb,mathtools,unicode-math}
\setmathfont{Latin Modern Math}
\usepackage{enumitem,needspace,xurl,hyperref}
\setlength{\parindent}{0pt}
\setlength{\parskip}{5pt}
\setlength{\emergencystretch}{3em}
\newcommand{\N}{\mathbb{N}}
\newcommand{\Z}{\mathbb{Z}}
\newcommand{\Q}{\mathbb{Q}}
\newcommand{\R}{\mathbb{R}}
\newcommand{\C}{\mathbb{C}}
\newcommand{\F}{\mathbb{F}}
\newcommand{\A}{\mathbb{A}}
\newcommand{\PP}{\mathbb{P}}
\newcommand{\E}{\mathbb{E}}
\newcommand{\eps}{\varepsilon}
\newcommand{\dd}{\,\mathrm{d}}
\begin{document}
\section*{UnsolvedMath ID: 30003762}
\textbf{Source problem code:} catalogue entry
\subsection{Definitions and mathematical statement}
For context, under the usual Galois-theoretic convention, a regular quadratic extension over a field \(K\) is a Galois extension \(E/K(T)\) of degree \(2\), where \(T\) is transcendental over \(K\) and \(K\) is algebraically closed in \(E\). Such an extension is called \(K\)-parametric if every quadratic Galois extension \(L/K\) is \(K\)-isomorphic to a specialization \(E_t/K\) at some \(t\in\mathbb{P}^{1}(K)\) outside the branch locus. These definitions give conventional context; the supplied record does not specify its base field or its precise convention for parametricity.

The supplied problem refers to a proportion \(f(H)\) of parametric regular quadratic extensions and records the estimate
\[
f(H)=1-O\!\left(\frac{\log H}{\sqrt H}\right)
\qquad(H\to\infty).
\]
If \(f(H)\) is a proportion, the displayed estimate means that there exist constants \(C>0\) and \(H_0>1\) such that, for every admissible height bound \(H\ge H_0\),
\[
0\le 1-f(H)\le C\frac{\log H}{\sqrt H}.
\]
The domain of admissible height bounds, the objects being counted, their equivalence relation, the height function, and the counting measure are not supplied, so this is only an interpretation of the recorded asymptotic notation, not a definition of its counting problem.

\textbf{Recorded question.} Can the error term in this estimate be removed unconditionally?

\textbf{Unresolved formulation.} The phrase ``removed unconditionally'' does not identify a precise mathematical conclusion. It could request an unconditional proof of a conditional estimate, a stronger error bound, or an exact equality such as \(f(H)=1\) for all sufficiently large admissible \(H\). The record does not choose among these assertions. Moreover, its dated literature assessment describes an unconditional gap involving scarcity of nontrivial rational points on quadratic twists, without specifying the curves, the meaning of nontriviality, or the implication connecting that scarcity to the displayed proportion. A precise conjecture cannot therefore be reconstructed from the supplied record alone.
\subsection{Short English statement}
The recorded problem asks whether an error term in a density estimate for parametric regular quadratic extensions can be removed unconditionally. The base field, counting family, height, and intended improvement are unspecified, so the exact assertion remains unresolved. The supplied open label and literature assessment require independent review.
\subsection{Sources}
\begin{itemize}
\item Supplied local UnsolvedMath registry record, permanent ID \(30003762\), exact source problem code \texttt{catalogue entry}, title: \emph{Sharp Density of Parametric Regular Quadratic Extensions}. The record gives proposed year 2018 and names no proposer.
\item Supplied source citation: \emph{Mini-Workshop: Arithmetic Geometry and Symmetries around Galois and Fundamental Groups} (2019), Oberwolfach Reports, pp.\ 1021--1080, DOI \texttt{10.4171/owr/2018/17}. Supplied source URL and literature reference ``OWR 2018 report'': \url{https://doi.org/10.4171/owr/2018/17}. The record describes these reports as documenting mathematical workshops and open-problem discussions.
\item Supplied literature reference: ``Granville, quadratic twists'', \url{https://doi.org/10.1093/imrn/rnm027}.
\item Supplied literature reference: ``Debes--Konig--Legrand--Neftin, parametric polynomials'', \url{https://arxiv.org/abs/2102.07465}.
\item Literature assessment dated exactly \texttt{2026-08-21}: this is supplied dated curation derived from the supplied OWR curation, not independently verified. It labels the problem open and states that Granville obtains the needed scarcity of nontrivial rational points on quadratic twists only conditionally on \(abc\), that later parametric-polynomial work does not remove the unconditional gap, and that the recorded general claim remains open. The supplied background explicitly requests independent verification of time-sensitive status and cited claims.
\item The registry labels the statement \texttt{corrected\_verified} and says that mathematical content was preserved while notation, typography, and omitted definitions were repaired for self-contained publication. This is supplied registry metadata, not independent mathematical verification of this draft.
\end{itemize}
\end{document}
